\documentclass[12pt,a4paper]{article}
\usepackage[utf8]{inputenc}
\usepackage[spanish]{babel}
\usepackage{graphicx}
\usepackage{listings}
\usepackage{xcolor}
\usepackage{hyperref}
\usepackage{amsmath}
\usepackage{geometry}
\usepackage{tikz}
\usepackage{float}
\usetikzlibrary{arrows.meta, positioning, shapes.geometric}
\geometry{margin=2.5cm}

% Configuración de listings para código MIPS
\lstdefinestyle{mipsstyle}{
    language={[MIPS]Assembler},
    basicstyle=\ttfamily\small,
    keywordstyle=\color{blue}\bfseries,
    commentstyle=\color{green!60!black},
    stringstyle=\color{red},
    numbers=left,
    numberstyle=\tiny,
    numbersep=5pt,
    frame=single,
    breaklines=true,
    showstringspaces=false,
    tabsize=4
}

\title{\textbf{Informe Técnico: Interrupciones y Excepciones en MIPS32} \\
\large Mecanismos de Gestión de Eventos Asíncronos y Síncronos}
\author{Arquitectura de Computadores}
\date{\today}

\begin{document}

\maketitle
\tableofcontents
\newpage

\section{Ciclo de una Interrupción de Hardware}

\subsection{Diagrama del Ciclo de Interrupción}

\begin{figure}[H]
\centering
\begin{tikzpicture}[node distance=1.5cm, auto, >=stealth']
    % Nodos del diagrama
    \node[rectangle, draw, fill=blue!20, minimum width=3cm, minimum height=1cm] (evento) {Evento Hardware};
    \node[rectangle, draw, fill=green!20, minimum width=3cm, minimum height=1cm, below=of evento] (irq) {Señal IRQ activada};
    \node[rectangle, draw, fill=yellow!20, minimum width=3cm, minimum height=1cm, below=of irq] (cpu) {CPU verifica interrupciones};
    \node[rectangle, draw, fill=orange!20, minimum width=3cm, minimum height=1cm, below=of cpu] (salto) {Salto a rutina de servicio};
    \node[rectangle, draw, fill=red!20, minimum width=3cm, minimum height=1cm, below=of salto] (contexto) {Guardar contexto};
    \node[rectangle, draw, fill=purple!20, minimum width=3cm, minimum height=1cm, below=of contexto] (servicio) {Ejecutar rutina ISR};
    \node[rectangle, draw, fill=red!20, minimum width=3cm, minimum height=1cm, below=of servicio] (restaurar) {Restaurar contexto};
    \node[rectangle, draw, fill=green!20, minimum width=3cm, minimum height=1cm, below=of restaurar] (retorno) {Retorno de interrupción};
    
    % Flechas
    \draw[->, thick] (evento) -- (irq);
    \draw[->, thick] (irq) -- (cpu);
    \draw[->, thick] (cpu) -- (salto);
    \draw[->, thick] (salto) -- (contexto);
    \draw[->, thick] (contexto) -- (servicio);
    \draw[->, thick] (servicio) -- (restaurar);
    \draw[->, thick] (restaurar) -- (retorno);
    
    % Descripciones laterales
    \node[right=of evento, xshift=1cm] (desc1) {Dispositivo genera interrupción};
    \node[right=of irq, xshift=1cm] (desc2) {Línea de interrupción activada};
    \node[right=of cpu, xshift=1cm] (desc3) {CPU verifica IE bit};
    \node[right=of salto, xshift=1cm] (desc4) {PC $\leftarrow$ Vector de interrupción};
    \node[right=of contexto, xshift=1cm] (desc5) {Guardar EPC, Status, registros};
    \node[right=of servicio, xshift=1cm] (desc6) {Atender al dispositivo};
    \node[right=of restaurar, xshift=1cm] (desc7) {Restaurar registros};
    \node[right=of retorno, xshift=1cm] (desc8) {eret $\rightarrow$ PC = EPC};
    
\end{tikzpicture}
\caption{Ciclo completo de una interrupción de hardware}
\end{figure}

\subsection{Descripción Paso a Paso}
\begin{enumerate}
    \item \textbf{Evento Hardware}: El dispositivo periférico completa una operación (ej: carácter recibido)
    \item \textbf{Activación IRQ}: El dispositivo activa su línea de interrupción
    \item \textbf{Verificación por CPU}: Al final de cada instrucción, la CPU verifica si hay interrupciones pendientes y si están habilitadas
    \item \textbf{Salto a rutina}: La CPU guarda la dirección de retorno en EPC y salta al vector de interrupción
    \item \textbf{Guardar contexto}: El software guarda los registros que va a utilizar
    \item \textbf{Ejecutar ISR}: Se ejecuta la rutina de servicio de interrupción
    \item \textbf{Restaurar contexto}: Se recuperan los registros guardados
    \item \textbf{Retorno}: Instrucción \texttt{eret} restaura el PC desde EPC
\end{enumerate}

\section{Sondeo vs. Interrupciones}

\subsection{Cuadro Comparativo}

\begin{table}[H]
\centering
\begin{tabular}{|p{5cm}|p{5cm}|}
\hline
\textbf{Sondeo (Polling)} & \textbf{Interrupciones} \\
\hline
CPU verifica constantemente el estado del dispositivo & Dispositivo notifica a la CPU cuando está listo \\
\hline
Espera activa (busy-waiting) & CPU puede ejecutar otras tareas mientras espera \\
\hline
Alta latencia de respuesta & Respuesta inmediata al evento \\
\hline
Simple de implementar & Mayor complejidad (hardware y software) \\
\hline
Ineficiente para eventos poco frecuentes & Eficiente para eventos asíncronos \\
\hline
Consume CPU constantemente & CPU ocupada solo durante la atención \\
\hline
\end{tabular}
\caption{Comparación entre sondeo e interrupciones}
\end{table}

\subsection{Ejemplo en MIPS32}
\begin{lstlisting}[style=mipsstyle, caption={Sondeo vs. Interrupción}]
# SONDEO: Espera activa
polling_teclado:
    lw $t0, KBD_CONTROL
    andi $t0, $t0, 1
    beqz $t0, polling_teclado   # Espera activa
    lw $v0, KBD_DATA            # Leer dato
    
# INTERRUPCIÓN: Configuración
configurar_int_teclado:
    li $t0, KBD_CONTROL
    li $t1, 0x3                 # Habilitar interrupción
    sw $t1, 0($t0)
    # CPU sigue con otras tareas
\end{lstlisting}

\section{Ventajas de las Interrupciones en el Uso del Procesador}

\begin{itemize}
    \item \textbf{Concurrencia}: Permite solapar operaciones de E/S con cómputo
    \item \textbf{Eficiencia energética}: CPU puede entrar en modos de bajo consumo mientras espera
    \item \textbf{Multitarea}: Base para sistemas operativos con tiempo compartido
    \item \textbf{Latencia}: Respuesta inmediata a eventos críticos
    \item \textbf{Rendimiento}: Aprovechamiento óptimo del CPU (hasta 100\% de utilización en cómputo útil)
\end{itemize}

\begin{figure}[H]
\centering
\begin{tikzpicture}
    \draw[->] (0,0) -- (10,0) node[right] {tiempo};
    
    % Sin interrupciones (polling)
    \foreach \x in {0,1.5,3,4.5,6,7.5,9} {
        \fill[red] (\x,0.5) rectangle ++(0.5,0.3);
    }
    \node at (5,1) {Polling: CPU siempre ocupada};
    
    % Con interrupciones
    \fill[green] (0,1.5) rectangle (2,1.8);
    \fill[green] (2.5,1.5) rectangle (4.5,1.8);
    \fill[green] (5,1.5) rectangle (7,1.8);
    \fill[green] (7.5,1.5) rectangle (9.5,1.8);
    
    \foreach \x in {2,4.5,7,9.5} {
        \fill[blue] (\x,1.5) rectangle ++(0.3,0.3);
    }
    \node at (5,2) {Interrupciones: CPU en cómputo útil (verde) con breves interrupciones (azul)};
    
\end{tikzpicture}
\caption{Comparativa de uso de CPU: Polling vs. Interrupciones}
\end{figure}

\section{Registros Especiales en MIPS32 para Interrupciones}

\begin{table}[H]
\centering
\begin{tabular}{|l|l|l|}
\hline
\textbf{Registro} & \textbf{Número} & \textbf{Función} \\
\hline
Status (SR) & $12 & Control de interrupciones (IE, IM, UM) \\
\hline
Cause & $13 & Identifica causa de excepción/interrupción \\
\hline
EPC & $14 & Dirección de retorno (PC interrumpido) \\
\hline
BadVAddr & $8 & Dirección que causó el fallo \\
\hline
Count & $9 & Temporizador (incrementa cada ciclo) \\
\hline
Compare & $11 & Valor de comparación para interrupción de timer \\
\hline
\end{tabular}
\caption{Registros CP0 relacionados con interrupciones}
\end{table}

\subsection{Bits importantes del registro Status}
\begin{itemize}
    \item \textbf{IEc (bit 0)}: Interrupt Enable global (0=deshabilitadas, 1=habilitadas)
    \item \textbf{IM (bits 8-15)}: Interrupt Mask - habilita interrupciones específicas
    \item \textbf{UM (bit 1)}: Modo usuario (1) o kernel (0)
    \item \textbf{EXL (bit 1)}: Exception Level (1 si estamos atendiendo excepción)
\end{itemize}

\begin{lstlisting}[style=mipsstyle, caption={Manipulación de registros CP0}]
# Leer registro Status
mfc0 $t0, $12

# Habilitar interrupciones globales
ori $t0, $t0, 0x1
mtc0 $t0, $12

# Deshabilitar interrupciones
mfc0 $t0, $12
andi $t0, $t0, 0xFFFFFFFE
mtc0 $t0, $12
\end{lstlisting}

\section{Necesidad de Guardar el Contexto}

\subsection{¿Por qué es necesario?}
\begin{enumerate}
    \item \textbf{Transparencia}: El programa interrumpido no debe notar la interrupción
    \item \textbf{Reanudación correcta}: Debe poder continuar exactamente donde se detuvo
    \item \textbf{Valores de registros}: La ISR puede modificar registros sin afectar al programa
    \item \textbf{Estado del procesador}: EPC y Status deben preservarse
\end{enumerate}

\subsection{Contexto mínimo a guardar}
\begin{itemize}
    \item Registros que modificará la ISR (por convención, $t0-$t9 si se usan)
    \item EPC (Exception Program Counter)
    \item Status (para restaurar estado de interrupciones)
\end{itemize}

\begin{lstlisting}[style=mipsstyle, caption={Guardado de contexto típico}]
ISR_teclado:
    # Guardar contexto
    sw $t0, contexto_t0
    sw $t1, contexto_t1
    sw $ra, contexto_ra
    mfc0 $t0, $14           # EPC
    sw $t0, contexto_epc
    
    # Atender interrupción
    # ... código de la ISR ...
    
    # Restaurar contexto
    lw $t0, contexto_t0
    lw $t1, contexto_t1
    lw $ra, contexto_ra
    lw $t0, contexto_epc
    mtc0 $t0, $14
    eret
\end{lstlisting}

\section{Momentos de Generación de Excepciones}

\subsection{a) Situaciones que generan excepciones}

\begin{enumerate}
    \item \textbf{Desbordamiento aritmético}: Resultado de operación excede rango
    \item \textbf{Fallo de dirección}: Acceso a memoria no alineado o fuera de rango
    \item \textbf{Instrucción ilegal}: Código de operación no válido
    \item \textbf{División por cero}: Operación aritmética inválida
    \item \textbf{Syscall}: Llamada al sistema explícita
    \item \textbf{Breakpoint}: Instrucción de depuración
    \item \textbf{Fallo de página}: Acceso a memoria no residente
    \item \textbf{Protección}: Acceso a región privilegiada desde modo usuario
\end{enumerate}

\subsection{b) Etapas del pipeline que provocan excepciones}

\begin{table}[H]
\centering
\begin{tabular}{|l|l|}
\hline
\textbf{Etapa} & \textbf{Excepciones posibles} \\
\hline
IF (Fetch) & Fallo de instrucción, dirección inválida \\
\hline
ID (Decode) & Instrucción ilegal, syscall, breakpoint \\
\hline
EX (Execute) & Desbordamiento aritmético, división por cero \\
\hline
MEM (Memory) & Fallo de dirección, fallo de página \\
\hline
WB (Write Back) & Generalmente no genera excepciones nuevas \\
\hline
\end{tabular}
\caption{Excepciones por etapa del pipeline}
\end{table}

\begin{figure}[H]
\centering
\begin{tikzpicture}
    \draw[->] (0,0) -- (12,0) node[right] {Ciclos};
    
    % Pipeline stages
    \foreach \x/\c in {0/IF, 1.5/ID, 3/EX, 4.5/MEM, 6/WB} {
        \draw[fill=blue!20] (\x,1) rectangle ++(1.2,0.5);
        \node at (\x+0.6,1.25) {\c};
    }
    
    % Excepción en EX
    \draw[red, thick, ->] (3.6,1.5) -- (3.6,2);
    \node[red] at (3.6,2.3) {Excepción aquí};
    
    \draw[red, dashed] (0,2.5) -- (12,2.5);
    \node at (1,2.7) {Instrucción anterior completa};
    \node at (9,2.7) {Instrucciones posteriores se descartan};
    
\end{tikzpicture}
\caption{Manejo de excepción en pipeline MIPS}
\end{figure}

\section{Estrategias de Tratamiento de Excepciones e Interrupciones}

\subsection{a) Diferencias entre Interrupciones y Excepciones}

\begin{table}[H]
\centering
\begin{tabular}{|l|l|}
\hline
\textbf{Interrupciones} & \textbf{Excepciones} \\
\hline
Asíncronas (eventos externos) & Síncronas (causadas por la instrucción) \\
\hline
Ocurren en cualquier momento & Ocurren en momentos predecibles \\
\hline
No están relacionadas con la instrucción actual & Causadas por la instrucción actual \\
\hline
Ej: teclado, timer, disco & Ej: división por cero, syscall \\
\hline
Se pueden enmascarar (deshabilitar) & Generalmente no enmascarables \\
\hline
\end{tabular}
\caption{Diferencias entre interrupciones y excepciones}
\end{table}

\subsection{b) Estrategias de Tratamiento en MIPS32}

\textbf{Estrategia 1: Vector de excepciones único}
\begin{itemize}
    \item Todas las excepciones saltan a la misma dirección (0x80000180)
    - El software debe leer el registro Cause para determinar la causa
\end{itemize}

\textbf{Estrategia 2: Vectores múltiples (MIPS32)}
\begin{itemize}
    \item Diferentes direcciones para diferentes tipos de excepción
    \item Configurable mediante registros CP0
\end{itemize}

\subsection{Redirección a la rutina de servicio}
\begin{enumerate}
    \item La hardware guarda PC en EPC
    \item Fuerza PC a la dirección del vector de excepciones
    \item Cambia a modo kernel (EXL bit en Status)
    \item Deshabilita interrupciones
\end{enumerate}

\subsection{Función del registro EPC}
\textbf{EPC (Exception Program Counter)} almacena la dirección de la instrucción que causó la excepción o la siguiente instrucción (según el tipo). Es fundamental para:
\begin{itemize}
    \item Reanudar la ejecución después de la excepción
    \item Instrucción \texttt{eret} copia EPC a PC
    \item Permite reintentar instrucciones (ej: fallo de página)
\end{itemize}

\section{Habilitación de Interrupciones}

\subsection{a) Habilitación en dispositivos y procesador}

\textbf{Teclado:}
\begin{lstlisting}[style=mipsstyle]
# Habilitar interrupciones en teclado
li $t0, KBD_CONTROL
li $t1, 0x3           # Bit0 = enable, Bit1 = interrupt enable
sw $t1, 0($t0)
\end{lstlisting}

\textbf{Pantalla:}
\begin{lstlisting}[style=mipsstyle]
# Habilitar interrupciones en pantalla
li $t0, CON_CONTROL
li $t1, 0x3           # Bit0 = ready, Bit1 = interrupt enable
sw $t1, 0($t0)
\end{lstlisting}

\textbf{Procesador (registro Status):}
\begin{lstlisting}[style=mipsstyle]
# Habilitar interrupciones globales
mfc0 $t0, $12
ori $t0, $t0, 0x1     # IEc = 1
mtc0 $t0, $12

# Habilitar interrupción específica (ej: timer)
mfc0 $t0, $12
ori $t0, $t0, 0x8000  # IM7 (timer) = 1
mtc0 $t0, $12
\end{lstlisting}

\subsection{b) Consecuencias de habilitar solo en dispositivos}
Si habilitamos interrupciones en dispositivos pero no en el procesador:
\begin{itemize}
    \item Los dispositivos generarán solicitudes de interrupción
    \item El procesador las ignorará (IEc = 0 en Status)
    \item Posibles consecuencias:
    \begin{itemize}
        \item Desbordamiento de buffers en dispositivos
        \item Pérdida de datos (caracteres de teclado no leídos)
        \item Dispositivos pueden quedar bloqueados esperando servicio
        \item Comportamiento indefinido del sistema
    \end{itemize}
\end{itemize}

\section{Procesamiento de Interrupciones}

\subsection{a) Interrupción de reloj paso a paso}

\textbf{Fase Hardware:}
\begin{enumerate}
    \item El temporizador (Count == Compare) genera señal de interrupción
    \item Se activa el bit correspondiente en registro Cause (IP)
    \item Al final de la instrucción actual, CPU verifica:
    \begin{itemize}
        \item IEc = 1 (interrupciones habilitadas)
        \item IM correspondiente = 1 (máscara habilitada)
        \item EXL = 0 (no estamos en excepción)
    \end{itemize}
    \item Guarda PC en EPC
    \item Guarda Status en un registro temporal
    \item Fuerza PC a 0x80000180 (vector de excepciones)
    \item Activa EXL = 1, deshabilita interrupciones
\end{enumerate}

\textbf{Fase Software (ISR):}
\begin{enumerate}
    \item Guarda contexto (registros generales en pila)
    \item Lee Cause para verificar que es interrupción de timer
    \item Reconoce la interrupción (resetea timer)
    \item Ejecuta rutina de servicio (ej: scheduler, contador de tiempo)
    \item Restaura contexto (registros generales)
    \item Ejecuta \texttt{eret} (restaura PC desde EPC, Status)
\end{enumerate}

\subsection{b) Importancia de guardar el contexto}
\begin{itemize}
    \item \textbf{Registros generales}: Si la ISR modifica registros, el programa interrumpido vería valores corruptos
    \item \textbf{EPC}: Necesario para reanudar la ejecución donde se detuvo
    \item \textbf{Status}: Para restaurar el estado de interrupciones previo
    \item \textbf{Anidamiento}: Permite que ocurran interrupciones dentro de la ISR si se re-habilitan
\end{itemize}

\section{Interrupciones de Reloj y Control de Ejecución}

\subsection{a) Usos de la interrupción de reloj}

\textbf{Prevenir bucles infinitos:}
\begin{lstlisting}[style=mipsstyle, caption={Timer como watchdog}]
ISR_timer:
    # Incrementar contador de ticks
    lw $t0, ticks
    addi $t0, $t0, 1
    sw $t0, ticks
    
    # Verificar tiempo máximo por programa
    lw $t1, tiempo_ejecucion
    bgt $t0, $t1, matar_programa
    
    # Reanudar
    j fin_ISR_timer
\end{lstlisting}

\textbf{Límite de tiempo de ejecución:}
\begin{itemize}
    \item Cada programa tiene un quantum de tiempo asignado
    \item La interrupción de reloj permite al scheduler cambiar de tarea
    \item Si un programa excede su tiempo, se fuerza un cambio de contexto
\end{itemize}

\subsection{b) Manejo si el programa termina antes}
\begin{enumerate}
    \item El programa termina voluntariamente (syscall exit)
    \item El sistema operativo remueve el programa de la cola de listos
    \item Se desactiva el timer asociado (o se ignora)
    \item Si ocurre la interrupción, la ISR detecta que no hay programa activo y simplemente retorna
\end{enumerate}

\begin{lstlisting}[style=mipsstyle]
ISR_timer:
    # Verificar si hay programa activo
    lw $t0, programa_activo
    beqz $t0, fin_ISR_timer   # No hay programa, ignorar
    
    # Manejar cambio de contexto normalmente
    # ...
\end{lstlisting}

\section{Análisis y Discusión de Resultados}

\subsection{Implementación de los Programas Prácticos}

\subsubsection{Buffer de Teclado con Filtro}
\begin{itemize}
    \item \textbf{Aciertos}: Uso correcto de buffer circular, filtrado eficiente de caracteres
    \item \textbf{Limitaciones}: Espera activa en lugar de interrupciones reales
    \item \textbf{Mejora}: Se podría implementar con interrupciones para eficiencia
\end{itemize}

\subsubsection{Semáforo con Pulsador}
\begin{itemize}
    \item \textbf{Máquina de estados}: Implementación clara y mantenible
    \item \textbf{Temporización}: Precisión adecuada para simulación
    \item \textbf{Interactividad}: Respuesta inmediata al pulsador
\end{itemize}

\subsection{Evaluación de los Mecanismos de Interrupción}

\textbf{Ventajas observadas:}
\begin{itemize}
    \item Capacidad de respuesta a eventos asíncronos
    \item Aislamiento entre tareas del sistema
    \item Escalabilidad para múltiples dispositivos
\end{itemize}

\textbf{Desafíos encontrados:}
\begin{itemize}
    \item Complejidad en la gestión de prioridades
    \item Necesidad de proteger secciones críticas
    \item Sobrecarga por cambio de contexto
\end{itemize}

\subsection{Conclusiones}

\begin{enumerate}
    \item Las interrupciones son fundamentales para sistemas interactivos y de tiempo real
    \item MIPS32 proporciona un mecanismo elegante y eficiente para gestión de excepciones
    \item El balance entre simplicidad (polling) y eficiencia (interrupciones) depende de la aplicación
    \item La correcta gestión del contexto es crítica para la estabilidad del sistema
    \item Los temporizadores son la base para sistemas multitarea por tiempo compartido
\end{enumerate}

\subsection{Recomendaciones para Implementaciones Reales}
\begin{itemize}
    \item Utilizar interrupciones en lugar de polling para E/S de baja frecuencia
    \item Implementar prioridades para interrupciones críticas
    \item Diseñar ISR cortas y eficientes
    \item Proteger secciones críticas con deshabilitado temporal de interrupciones
    \item Utilizar el timer para watchdog y time-sharing
\end{itemize}

\end{document}